\article{Drinfeld 的 shtuka 导论}{An introduction to Drinfeld's ``Shtuka''}{D. A. Kazhdan}{2}{347--356}

先从特征为 $p>0$ 的局部域（local field）$K$ 上 $\GL(n)$ 的 Langlands 猜想说起。固定素数 $l\ne p$，用 $\overline\Q_l$ 表示 $\Q_l$ 的一个代数闭包（algebraic closure）。\footnote{AMS（MOS）主题分类（1970）：主要 10H10；次要 10D20、22E55、20G35。}

先约定记号。设 $G$ 为拓扑群（topological group），$U\subset G$ 为紧子群（compact subgroup）。考虑 $G$ 的光滑（smooth）$\overline\Q_l$ 表示，即同态 $\rho:G\longrightarrow\Aut V$，其中 $V$ 为 $\overline\Q_l$ 向量空间，且每个向量 $v\in V$ 的稳定子群（stabilizer）都是开的。若对每个紧开子群（compact open subgroup）$U\subset G$ 都有 $\dim V^U<\infty$，则称 $\rho$ 是容许的（admissible），其中 $V^U$ 为 $U$ 不变向量的集合。

\paragraph*{定义（Definition）.}
用 $\widehat G$ 表示 $G$ 的不可约（irreducible）容许表示的等价类（equivalence class）集合，用 $\widehat G_n\subset\widehat G$（$n=1,2,\ldots$）表示由 $n$ 维表示的等价类构成的子集。

\paragraph*{注（Remark）.}
容易验证，容许 $\C$ 表示的标准结果（见[7]）对这里的 $\overline\Q_l$ 表示同样成立。

现在可以表述局部非分歧互反律（local unramified reciprocity law）。设 $K_{\mathfrak p}$ 为非 Archimedes 局部域（nonarchimedean local field），$o=o_{\mathfrak p}$ 为其整数环（ring of integers），$\pi\in o$ 为素元（prime element），$k=o/(\pi)$ 为剩余域（residue field），$q=\operatorname{card}k$，$v$ 为满足 $v(\pi)=1$ 的 $K_{\mathfrak p}$ 上的赋值（valuation）。

回忆，$\GL(n,K_{\mathfrak p})$ 的不可约容许表示若含有非零的 $\GL(n,o)$ 不变向量，就称为非分歧（unramified）。

\paragraph*{注.} 此时这一维数为一。

$\GL(n,K_{\mathfrak p})$ 的非分歧表示等价类集合 $\widehat G_{\mathfrak p,\mathrm{un}}$ 的分类是熟知的。设 $B\subset\GL(n,K_{\mathfrak p})$ 为三角矩阵子群。对每个 $x=(x_1,\ldots,x_n)$（$x_i\in\overline\Q_l^*$），定义 $B$ 的特征（character）$\chi_x$：
\[
 \chi_x(b)=(q^{1-n}x_1)^{v(b_{11})}(q^{2-n}x_2)^{v(b_{22})}\cdots x_n^{v(b_{nn})},
\]
其中 $b_{ii}$ 为 $b$ 的对角元。用 $\overline\rho_x$ 表示诱导表示（induced representation）$\operatorname{ind}_B^{\GL(n,K_{\mathfrak p})}\chi_x$；它实现在满足下式的 $\GL(n,K_{\mathfrak p})$ 上的局部常值函数（locally constant function）空间中：
\[
 f(bg)=\chi_x(b)f(g),\qquad b\in B,\quad g\in\GL(n,K_{\mathfrak p}).
\]
表示 $\overline\rho_x$ 未必不可约，但具有有限的 Jordan--Hölder 列（Jordan--Hölder series），其中恰有一个非分歧组成因子。记此因子为 $\rho_x$。下面的断言是熟知的[1]。

\begin{originaltheorem}{定理 1（Theorem 1）}
(a) 对称群（symmetric group）$S_n$ 的每个元素 $s$ 都满足 $\rho_x\simeq\rho_{s(x)}$。

(b) 映射
\[
 \rho:(\overline\Q_l^*)^n/S_n\longrightarrow\widehat G_{\mathfrak p,\mathrm{un}}
\]
是同构（isomorphism）。
\end{originaltheorem}

% 原书348页
\paragraph*{注.} 这里采用的参数化与通常的参数化相差一个平移，从而避免了 $q$ 的分数次幂。

用 $\overline K_{\mathfrak p}$ 表示 $K_{\mathfrak p}$ 的一个可分闭包（separable closure），$\overline k_{\mathfrak p}$ 表示相应的 $k_{\mathfrak p}$ 的代数闭包。令 $W_{K_{\mathfrak p}}$ 为 $\Gal(\overline K_{\mathfrak p}/K_{\mathfrak p})$ 中由那些在 $\overline k_{\mathfrak p}$ 上诱导映射 $x\mapsto x^{q^n}$（某个 $n\in\Z$）的元素构成的稠密子群（dense subgroup），用 $I_{K_{\mathfrak p}}\subset W_{K_{\mathfrak p}}$ 表示固定 $\overline k_{\mathfrak p}$ 的子群。因此 $W_{K_{\mathfrak p}}$ 为 $K_{\mathfrak p}$ 的 Weil 群（Weil group），$I_{\mathfrak p}$ 为惯性群（inertia group）。

若 $W_{K_{\mathfrak p}}$ 的一个表示完全可约（completely reducible），且其在 $I_{K_{\mathfrak p}}$ 上的限制平凡，则称该表示非分歧。由此便可表述非分歧表示的局部互反律，其证明也很直接。

\begin{originaltheorem}{定理--定义 2（Theorem--Definition 2）}
$\widehat G_{\mathfrak p,\mathrm{un}}$ 与 $W_{K_{\mathfrak p}}$ 的 $n$ 维非分歧表示的等价类集合之间存在自然的一一对应。
\end{originaltheorem}
\paragraph*{注.} 这里 $n$ 维表示指的是到 $\GL(n,\overline\Q_l)$ 的态射（morphism）。
\begin{proof}[证明与构造（Proof--Construction）]
设 $\tau$ 为 $W_{K_{\mathfrak p}}$ 的非分歧半单表示（semisimple representation）。由于它非分歧，可以视其为
$W_{K_{\mathfrak p}}/I_{K_{\mathfrak p}}\simeq\Z$ 的表示。因此它对应于 $\GL(n,\overline\Q_l)$ 的一个元素 $\alpha_\tau$，即 Frobenius 元素 $\operatorname{Fr}$ 的像。因为 $\tau$ 半单，$\alpha_\tau$ 也半单，所以 $\alpha_\tau$ 的共轭类（conjugacy class）完全由其特征值（eigenvalue）集合 $x=(x_1,\ldots,x_n)\in(\overline\Q_l^*)^n$ 决定。令 $\tau$ 对应于 $\GL(n,K_{\mathfrak p})$ 的表示 $\rho=\rho_x$，记作 $\rho\sim\tau$。
\end{proof}

这一构造使用了非分歧表示的显式实现，但还存在更为典范的表述。已知（见[1]），对于每个自然数 $m$，$\GL(n,K_{\mathfrak p})$ 上存在唯一的 $\Q$ 值紧支集（compact support）局部常值函数 $\chi_m$，满足：(a) $\chi_m$ 为双 $\GL(n,o)$ 不变；(b) 对于每个 $x=(x_1,\ldots,x_n)\in(\overline\Q_l^*)^n$，
\[
 \Tr\rho_x(\chi_m)=x_1^m+\cdots+x_n^m.
\]
现在可以将定理2重新表述如下。

\begin{originaltheorem}{定理 $2'$（Theorem $2'$）}
$\GL(n,K_{\mathfrak p})$ 的不可约非分歧表示 $\rho$ 对应于 $W_{K_{\mathfrak p}}$ 的 $n$ 维非分歧表示 $\tau$，当且仅当对所有自然数 $m>0$，
\[
 \Tr\rho(\chi_m)=\Tr\tau(\operatorname{Fr}^m).
\]
\end{originaltheorem}

现在考虑整体情形。设 $K$ 为整体域（global field），$\operatorname{char}K=p$。可以把 $K$ 视为有限域（finite field）$k$ 上光滑绝对不可约曲线（smooth absolutely irreducible curve）$X$ 的有理函数域。令 $q=\operatorname{card}k$，用 $\Pi$ 表示 $K$ 的全体素点。对每个 $\mathfrak p\in\Pi$，用 $K_{\mathfrak p}$ 表示 $K$ 在 $\mathfrak p$ 处的完备化（completion）。用 $\A$ 表示 $K$ 的阿代尔环（adele ring），并考虑 $\overline\Q_l$ 空间 $L$，其元素为
$\GL(n,K)\backslash\GL(n,\A)$ 上取 $\overline\Q_l$ 值的局部常值尖点函数（cuspidal function）。这就是说，对于每个 $K$ 子空间 $\Lambda\subset K^n$，$\Lambda\ne0,K^n$，有
\[
 \int_{U_\Lambda(K)\backslash U_\Lambda(\A)}f(ux)\,du=0,
\]
其中 $U_\Lambda$ 为满足 $g\Lambda=\Lambda$，且在 $\Lambda$ 及 $K^n/\Lambda$ 上作用平凡的 $g\in\GL(n,K^n)$ 构成的子群。$\GL(n,\A)$ 以右平移作用于 $L$。每个 $f\in L$ 的支集模去 $\A^*$（即 $\GL(n,\A)$ 的中心）后是紧的（例如见 Godement--Jacquet, Lecture Notes in Math., vol. 260, Springer, 第142页）；因此 $\GL(n,\A)$ 在 $L$ 上的表示 $T$ 是光滑的。

% 原书349页
若 $\GL(n,\A)$ 的一个不可约容许表示作为 $L$ 的子表示（subrepresentation）出现，则称其为尖点自守表示（cuspidal automorphic representation）。用 $\widehat{\GL(n,\A)}_c$ 表示这些表示的集合。

由[5]等文献可知，$\GL(n,\A)$ 的任意不可约容许表示 $(\rho,V)$ 均可唯一表示为局部不可约容许表示的限制张量积（restricted tensor product）$\bigotimes_{\mathfrak p\in\Pi}\rho_{\mathfrak p}$。若 $S$ 是 $\Pi$ 的有限子集，且对 $\mathfrak p\notin S$，$\rho_{\mathfrak p}$ 非分歧，则称 $\rho$ 在 $S$ 外非分歧。

下述断言（见[10]）有助于精确表述 Langlands 猜想。
\begin{originaltheorem}{定理 3（Theorem 3）}
设 $\rho,\rho'$ 为 $\GL(n,\A)$ 的两个尖点自守表示，且对几乎所有 $\mathfrak p\in\Pi$，有 $\rho_{\mathfrak p}=\rho'_{\mathfrak p}$。则 $\rho=\rho'$。
\end{originaltheorem}

现在考虑 Galois 一侧。设 $\overline K$ 为 $K$ 的可分代数闭包，$W_K$ 为 $\Gal(\overline K/K)$ 中在 $\overline k$ 上诱导 $x\mapsto x^{q^n}$（某个 $n\in\Z$）的元素构成的稠密子群。对于每个 $\mathfrak p\in\Pi$，有一个嵌入（embedding）$W_{K_{\mathfrak p}}\hookrightarrow W_K$，确定到 $W_K$ 中的共轭（见[12]）。

\paragraph*{定义.} 设 $\tau$ 为 $W_K$ 的有限维 $\overline\Q_l$ 表示。若 $\tau$ 在 $W_{K_{\mathfrak p}}$ 上的限制 $\tau_{\mathfrak p}$ 非分歧，则称 $\tau$ 在 $\mathfrak p$ 处非分歧。

\paragraph*{注 1.} 已知对每个这样的 $\tau$，存在有限集合 $S\subset\Pi$，使得 $\tau$ 在 $S$ 外非分歧。
\paragraph*{注 2.} 若 $\tau$ 在 $\mathfrak p$ 处非分歧，则 $\tau(\operatorname{Fr}_{\mathfrak p})$ 确定到共轭。

\begin{originaltheorem}{定理 4（Theorem 4，见[12]）}
设 $\tau,\tau'$ 为 $W_K$ 的两个不可约表示，且对几乎所有 $\mathfrak p\in\Pi$，$\tau(\operatorname{Fr}_{\mathfrak p})$ 与 $\tau'(\operatorname{Fr}_{\mathfrak p})$ 共轭。则 $\tau$ 与 $\tau'$ 等价。
\end{originaltheorem}

现在可以表述 \emph{Langlands 猜想}：集合 $\widehat{\GL(n,\A)}_c$ 与集合 $(\widehat W_K)_n$ 之间存在一一对应 $\varphi:\rho\leftrightarrow\tau$，并且对每个尖点自守表示 $\rho$，存在有限集合 $S\subset\Pi$，使得 $\rho$ 和 $\tau$ 均在 $S$ 外非分歧，且对所有 $\mathfrak p\in\Pi-S$，有 $\rho_{\mathfrak p}\sim\tau_{\mathfrak p}$（见定理2）。

\paragraph*{注 1.} 定理3和定理4蕴含这种对应（若存在）的唯一性。
\paragraph*{注 2.} 还有另一种猜想，要求尖点自守表示与 $W_K$ 的 $l$ 进表示（$l$-adic representation）系统之间有一一对应。由 Deligne [3] 和 Drinfeld 的结果可知，此猜想在 $n=2$ 时成立。

\paragraph*{注.} 由整体猜想可推出下述局部猜想。设 $M$ 为正特征局部域。$\GL(n,M)$ 的超尖点表示（supercuspidal representation）$\rho_M$ 与 $W_M$ 的 $n$ 维不可约表示 $\tau_M$ 之间存在一一对应。该对应满足：对于任意整体域 $K$、任意满足 $M\simeq K_{\mathfrak p_0}$ 的点 $\mathfrak p_0\in\Pi$，以及任意满足 $\rho_{\mathfrak p_0}$ 超尖点的自守表示 $\rho=\bigotimes\rho_{\mathfrak p}$，$W_K$ 的表示 $\tau\sim\rho$ 在 $W_{K_{\mathfrak p_0}}\simeq W_M$ 上的限制 $\tau_{\mathfrak p_0}$ 对应于 $\rho_{\mathfrak p_0}$。

这引出一个更一般的问题：如何在两个群（例如 $G$ 和 $W$）的表示之间建立一一对应？

% 原书350页
一种方法是，对 $G$ 的每个表示 $\rho$，在某个集合 $X$ 上定义函数 $\Gamma_\rho$，同样，对 $W$ 的每个表示 $\tau$ 定义函数 $\Gamma_\tau$。然后规定，$\rho\sim\tau$ 当且仅当 $\Gamma_\rho=\Gamma_\tau$。当 $G=\GL(2,\A)$ 或 $\GL(3,\A)$、$W=W_K$、$X$ 为伊代尔类群（idele class group）$K^*\backslash\A^*$ 的特征集合，而 $\Gamma_\rho$ 和 $\Gamma_\tau$ 为 $\Gamma$ 函数时，这种方法非常有效（见[7]、[9]），由此，对 $W_K$ 的每个二维或三维不可约 $l$ 进表示系统 $\tau$，都能找到相应的 $\GL(2,\A)$ 或 $\GL(3,\A)$ 的表示 $\rho$。然而迄今为止，这种途径还不能使我们反过来从 $\GL(n,\A)$ 的表示 $\rho$ 构造 $W_K$ 的表示 $\tau$。

H. Weyl [13] 最早使用了另一种更为显式的方法，在两个群 $G,W$ 的表示之间建立对应。他考虑的是 $G=\GL(n,\R)$、$W=S_m$ 为对称群的情形。他研究 $G\times W$ 在 $V\simeq\R^m$ 的张量幂（tensor power）$L=V^{\otimes m}$ 上的表示。\footnote{\textbf{译者注（本次校订）：}这里应为 $V\simeq\R^n$。$\GL(n,\R)$ 通过其 $n$ 维标准表示作用于每个张量因子，$S_m$ 则置换 $m$ 个因子；下文的条件 $n\ge m$ 也与此一致。}将 $L$ 分解为 $G\times W$ 不变不可约子空间的直和（direct sum）$L=\bigoplus L_i$。众所周知，$L_i$ 为张量积 $L_i=\rho_i\otimes\tau_i$，其中 $\rho_i$ 为 $\GL(n,\R)$ 的不可约表示，$\tau_i$ 为 $S_m$ 的不可约表示。他发现，当 $n\geq m$ 时，以这种方式得到了 $S_m$ 的表示与 $\GL(n,\R)$ 的部分表示之间的一一对应。这意味着：(a) $\{\tau_i\}=\widehat S_m$；(b) 若 $\rho_i\simeq\rho_j$，则 $i=j$。

Drinfeld 把这种方法用于 $G=\GL(2,\A)$、$W=W_K$ 的情形。最初的想法是构造 $\GL(2,\A)\times W_K$ 的一个 $\overline\Q_l$ 表示 $(T,L)$\footnote{由于 $(T,L)$ 由这些性质唯一确定到同构，自然可以期待它定义在 $\Q$ 上，而不只是在 $\overline\Q_l$ 上。}，使其具有下述性质：考虑 $L$ 分解成 $G\times W$ 不变不可约子空间的和 $L=\bigoplus L_i$，并将每个 $L_i$ 写作 $\rho_i\otimes\tau_i$，其中 $\rho_i\in\widehat{\GL(2,\A)}_c$、$\tau_i\in(\widehat W_K)_2$，则 (a) $\{\rho_i\}=\widehat{\GL(2,\A)}_c$；(b) 对所有 $i$，$\tau_i\sim\rho_i$。但遗憾的是，这样的表示 $(T,L)$ 不存在。

为说明这一障碍，先回顾几个定义和结果。它们在有限群表示论中是标准的（见[11]），也容易推广到容许表示。

设 $G$ 为群，$(\rho,V)$ 为 $G$ 的一个不可约容许 $\overline\Q_l$ 表示。用 $M(\rho)\subset\overline\Q_l$ 表示 $\rho$ 的等价类的定义域（field of definition）。具体地，$M(\rho)$ 是以下子群的不动域： $g_\rho\subset\Gal(\overline\Q_l/\Q_l)$，后者由满足 $\rho^\sigma\simeq\rho$ 的元素 $\sigma\in\Aut(\overline\Q_l:\Q_l)$ 构成。

\paragraph*{定义.} 若 $\rho$ 能在 $M(\rho)$ 上实现，则称 $\rho$ 无障碍（unobstructed）。这就是说，存在 $M(\rho)$ 上的向量空间 $V_0$ 与态射 $\rho_0:G\longrightarrow\Aut V_0$，使得 $(\rho_0,V_0\otimes_{M(\rho)}\overline\Q_l)$ 与 $(\rho,V)$ 等价。

下面两个断言是熟知的（见[11]）。
\begin{originaltheorem}{引理 1（Lemma 1）}
设 $(\rho,V)$ 为 $G$ 的表示，$H\subset G$ 为子群，且 $\dim V^H=1$。则 $\rho$ 无障碍。
\end{originaltheorem}
\begin{originaltheorem}{引理 2（Lemma 2）}
设 $L$ 是 $G\times W$ 的光滑 $\overline\Q_l$ 表示的空间，$(\rho,V)\in\widehat G$、$(\tau,C)\in\widehat W$，且作为 $G$ 模有 $\Hom_W(C,L)\simeq V$。若 $\rho$ 无障碍，则 $\tau$ 也无障碍。\footnote{\textbf{译者注（本次校订）：}按此处写出的条件，引理还缺少关于 $L$ 的有理定义及下降相容性的假设。取 $G$ 为平凡群、$L=C$ 为任意不可约表示 $\tau$ 的空间，则 $\Hom_W(C,L)$ 是一维平凡表示；若不补充条件，引理就会断言一切 $\tau$ 都无障碍，与下文所举情形矛盾。因此，前后文的“不存在”应针对具有所需有理下降结构的实现；单就 $\overline\Q_l$ 表示而言，给定对应后总可以取所写的抽象直和。}
\end{originaltheorem}
\begin{originaltheorem}{引理 3（Lemma 3）}
$\GL(2,\A)$ 的每个不可约表示都无障碍。
\end{originaltheorem}

% 原书351页
\begin{proof}
对于 $K$ 的每个整理想（positive ideal）$\mathfrak a$，用 $\Gamma_{\mathfrak a}\subset\GL(2,\A)$ 表示由整元素矩阵构成、模 $\mathfrak a$ 保持向量 $(0,1)$ 不变的子群。已知（见[10]）存在理想 $\mathfrak a$，使得 $\dim V^{\Gamma_{\mathfrak a}}=1$。于是由引理1得到结论。

另一方面，容易构造一个有障碍的 $W_K$ 的二维表示 $\tau$。例如，可令 $\tau$ 由 $W_K$ 中某个指数为二的子群的一维表示诱导而来。此时可以构造（见[1]）$\rho\in\widehat{\GL(2,\A)}_c$，使得 $\rho\sim\tau$。于是可见，具有性质(a)、(b)的表示 $(T,L)$ 不存在。
\end{proof}

可以如何处理这一障碍？对任意群 $H$ 的不可约表示 $\tau$，$H\times H$ 的表示 $\tau\otimes\widehat\tau$ 都无障碍，其中 $\widehat\tau$ 为 $\tau$ 的逆反表示（contragredient representation）。所以可以尝试实现
\[
 \bigoplus\rho\otimes\tau\otimes\widehat\tau,
 \qquad \rho\in\widehat{\GL(2,\A)}_c,\quad\tau\sim\rho.
\]
这一构想已接近可行。进一步实现它时，还需要引入直积分（direct integral）。

为精确表述 Drinfeld 的结果，还需要一些定义。
\paragraph*{定义.} 设 $(\rho,V)$ 为 $\GL(n,\A)$ 的一个表示。若可将 $V$ 写为直和 $V=\bigoplus_{t\in\Z}V^t$，使得对每个 $g\in\GL(n,\A)$，
\[
 \rho(g)V^t=V^{t+v(\det(g))},
\]
则称 $\rho$ 为分次表示（graded representation）。

\paragraph*{注.} 显然，对于任意分次表示 $(\rho,V)$，算子 $\rho(f)$（$f\in C_c^\infty(\GL(n,\A))$）不是迹类（trace class），无法定义特征（character）$\Tr_\rho$。但有时可以定义正规化特征（regularized character）
\[
 \operatorname{Tr}^{\mathrm{reg}}_\rho(f)
 \mathrel{\mathop:}=\Tr P_0\rho(f)P_0,
 \qquad f\in C_c^\infty(\GL(n,\A)),
\]
其中 $P_0$ 为到 $V^0$ 的投影。容易证明，$\operatorname{Tr}^{\mathrm{reg}}_\rho$ 的定义与分次 $V=\bigoplus V^t$ 的选取无关。\footnote{译注：原文用特殊叠印的迹符号表示正规化迹；这里统一记为 $\operatorname{Tr}^{\mathrm{reg}}$，以区别通常的迹 $\Tr$。}

设 $(\rho,V)$ 为分次表示，$\chi\in\widehat{\A^*}$。以 $V_\chi$ 表示 $V$ 对下列向量所张成的子空间取商所得的空间：\footnote{\textbf{译者注（本次校订）：}此处应当按中心特征取商。若按下式让 $g$ 遍历整个 $\GL(n,\A)$，所得商表示上必有 $\rho(g)=\chi(\det g)$，因而只能含一维组成因子，与定理6中的尖点表示不符。通常的定义是模去 $\rho(zI_n)v-\chi(z)v$（$z\in\A^*,v\in V$）张成的子空间，使中心通过 $\chi$ 作用。}
\[
 \{\chi(\det g)v-\rho(g)v:v\in V,\ g\in\GL(n,\A)\}.
\]
\paragraph*{定义.} 若对于所有 $\chi\in\widehat{\A^*}$，$V_\chi$ 都是不可约表示的直和，则称 $(\rho,V)$ 完全可约。

下面的结果容易证明。
\begin{originaltheorem}{引理（Lemma）}
设 $(\rho,V)$、$(\rho',V')$ 是两个分次完全可约表示，且 $\operatorname{Tr}^{\mathrm{reg}}_\rho$、$\operatorname{Tr}^{\mathrm{reg}}_{\rho'}$ 都存在并相等。则 $\rho\simeq\rho'$。
\end{originaltheorem}

现在可以表述 Drinfeld 结果的“构造性”形式。
\begin{originaltheorem}{定理 6（Theorem 6）}
存在 $\GL(2,\A)\times W_K\times W_K$ 的表示 $(\rho,V)$，满足：

(a) $\rho$ 在 $K^*\subset\operatorname{Center}\GL(2,\A)$ 上的限制平凡；

(b) 对于每个 $\chi\in\widehat{\A^*/K^*}$，$V_\chi$ 是不可约表示的直和
\[
 V_\chi=\bigoplus_i V_{\chi,i},\qquad
 V_{\chi,i}=\rho_{\chi,i}\otimes\tau_{\chi,i}\otimes\widehat\tau_{\chi,i};
\]

(c) $\widehat{\GL(2,\A)}_c=\{\rho_{\chi,i}\}_{\chi,i}$；

(d) 对所有 $(\chi,i)$，$\rho_{\chi,i}\sim\tau_{\chi,i}$。
\end{originaltheorem}

% 原书352页
这个定理的证明分为两步：(A) 构造 $(\rho,V)$；(B) 证明 $(\rho,V)$ 满足条件(a)--(d)。

我们只简要讨论这两部分。首先考虑 $W_K\times W_K$ 的几何解释。设 $\Lambda$ 为 $X\times X$ 的函数域，$\Lambda_s\supset\Lambda$ 为其可分闭包，$W_\Lambda\subset\Gal(\Lambda_s/\Lambda)$ 为在 $\overline k$ 上诱导 $x\mapsto x^{q^n}$（$n\in\Z$）的元素构成的子群。

有自然映射 $\pi:W_\Lambda\longrightarrow W_K\times W_K$，它就是将任意 $\sigma\in W_\Lambda\subset\Aut(\Lambda_s/\Lambda)$ 限制到 $K_s\otimes K_s\subset\Lambda_s$。显然，$\operatorname{Im}\pi$ 由满足 $\sigma'|\overline k=\sigma''|\overline k$ 的点对 $(\sigma',\sigma'')\in W_K\times W_K$ 构成。Drinfeld 定义了扩张（extension）
\[
 0\longrightarrow W_\Lambda\longrightarrow\widetilde W_\Lambda
 \longrightarrow\Z\longrightarrow0\tag{**}
\]
及满同态（epimorphism）$\widetilde\pi:\widetilde W_\Lambda\longrightarrow W_K\times W_K$。

为此，考虑 $\Lambda$ 的代数闭包 $\overline\Lambda$ 的自同构 $\gamma$ 构成的群 $\dot W_\Lambda$，要求 $\gamma$ 在完美闭包（perfect closure）$\Lambda^{K_{\mathrm{per}}\otimes K_{\mathrm{per}}}_{\mathrm{per}}\subset\overline\Lambda$ 上的限制具有 $\operatorname{Fr}_1^m\operatorname{Fr}_2^n$ 的形式，其中 $\operatorname{Fr}_1,\operatorname{Fr}_2$ 为 $K_{\mathrm{per}}\otimes K_{\mathrm{per}}$ 上的部分 Frobenius（partial Frobenius）。定义 $\widetilde W$ 为 $\dot W_\Lambda$ 模去由 Frobenius 生成的子群所得的商群。

有自然嵌入 $W_\Lambda\longrightarrow\dot W_\Lambda$，它给出扩张(**)。
\paragraph*{注.} 类似地可以定义群 $\dot W_K$ 和 $\widetilde W_K$，但此时 $\widetilde W_K=W_K$。

态射 $\pi$ 扩张为
\[
 \widetilde\pi:\widetilde W_\Lambda\longrightarrow
 \widetilde W_K\times\widetilde W_K=W_K\times W_K.
\]
\begin{originaltheorem}{定理 7（Theorem 7，Drinfeld）}
(1) $\widetilde\pi$ 为满同态。

(2) 对于每个有限群 $H$ 及每个同态 $\varphi:\widetilde W_\Lambda\longrightarrow H$，可将 $\varphi$ 写为复合 $\varphi=\overline\varphi\circ\widetilde\pi$，其中 $\overline\varphi:W_K\times W_K\longrightarrow H$ 为同态。
\end{originaltheorem}

这一定理的证明不难，但完整写出证明需要另引入若干定义，因此留给读者。

接下来考察如何构造表示 $(\rho,V)$。已知可以将 $K$ 视为有限域 $k$ 上光滑、射影、绝对不可约曲线 $X$ 的函数域。令 $S=X\times X$。假设可以在 $S$ 的泛点（generic point）$\eta$ 上定义代数簇（algebraic variety）的一个投射系统（projective system）$M_i$（$i\in I$），满足：

(a) $\GL(2,\A)$ 作用于投射极限（projective limit）$M=\varprojlim M_i$（把所有 $M_i$ 看成 $\Lambda$ 簇）；

(b) 对每个 $i\in I$，部分 Frobenius $\operatorname{Fr}_1,\operatorname{Fr}_2$（看作 $k(\eta)$ 的自同态）都有到 $M_i$ 的自同态的提升，且 $\operatorname{Fr}_1\circ\operatorname{Fr}_2$ 提升为 $M_i$ 上的 Frobenius。

于是可以构造 $\GL(2,\A)\times W_K\times W_K$ 的表示 $(\rho_M,V_M)$。为此，固定任意整数 $r$，考虑
\[
 V_M=\varinjlim H_c^r(M_i,\overline\Q_l).\footnotemark
\]
\footnotetext{$H_c^r$ 表示紧支集 $l$ 进上同调（$l$-adic cohomology with compact support）。}
群 $\GL(2,\A)\times\widetilde W_\Lambda$ 自然作用于 $V_M$，且对每个 $i\in I$，$\widetilde W_\Lambda$ 保持 $H_c^r(M_i,\overline\Q_l)$ 在 $V_M$ 中的像。由定理7，该作用可经由 $\GL(2,\A)\times W_K\times W_K$ 分解。

Drinfeld 的构造实际上遵循这个方案。但簇 $M_i$ 不是有限型（finite type）。这一点可从构造的目标看出。我们希望得到 $\GL(2,\A)\times W_K\times W_K$ 的一个表示 $(\rho,V)$，使其限制到 $\GL(2,\A)$ 后，为 $\GL(2,K)\backslash\GL(2,\A)$ 上局部常值尖点函数空间 $L$ 中的表示的四倍。按定义，$L$ 嵌入一个更大的空间 $\check L$，即 $\GL(2,K)\backslash\GL(2,\A)$ 上全体紧支集局部常值函数的空间。因此可以期待，实现 $(\rho,V)$ 的唯一方法，是将其作为 $(\check\rho,\check V)$ 的子空间实现，且 $(\check\rho,\check V)$ 限制到 $\GL(2,\A)$ 后是 $\check L$ 的某个倍数。可以尝试将 $\check V$ 实现为 $\varprojlim H_c^r(M_i,\overline\Q_l)$，再将 $V$ 实现为 $\check V$ 的子空间。

% 原书353页
但表示 $\check L$ 是“大的”。这意味着，对于紧开子群 $U\subset\GL(2,\A)$ 及特征 $\chi\in\widehat{\A^*/K^*}$，$\check L_\chi$ 中 $U$ 不变向量的子空间 $(\check L_\chi)^U$ 为无限维。这意味着（见下文）$\dim H^r(M_i^0,\overline\Q_l)=\infty$\footnote{$M_i^0$ 为 $M_i$ 的一个连通分支。}，从而 $M_i^0$ 不可能是有限型簇。

以上说明了构造中必然出现的复杂性。下面描述 Drinfeld 构造的主要结构，略去具体定义。

对 $X$ 上每个有效除子（positive divisor）$D$，Drinfeld 在 $\eta$ 上定义一个代数空间（algebraic space）$M_D$；对每个 $D'\supset D$，定义态射 $M_{D'}\longrightarrow M_D$。这个空间是连通不可约分支的并
\[
 M_D=\bigcup_{n=-\infty}^\infty M_D^n.
\]
所有偶数次分支 $M_D^{2n}$ 彼此同构，所有奇数次分支 $M_D^{2n+1}$ 也彼此同构，其中 $n\in\Z$。它们具有下述结构：(1) 每个 $M_D^n$ 是有限型曲面（surface）$M_D^{K,n-K}$ 的并
\[
 M_D^n=\bigcup_K M_D^{K,n-K};
\]
(2) $M_D^{K+1,n-K-1}\supset M_D^{K,n-K}$；(3) 差集 $M_D^{K+1,n-K-1}-M_D^{K,n-K}$ 是仿射直线（affine line）的并。Drinfeld 还为每个子空间 $M_D^{K,n-K}$ 定义紧化（compactification）$\overline M_D^{K,n-K}$，并证明双有理同构（birational isomorphism）$M_D^{K+1,n-K-1}\sim M_D^{K,n-K}$ 可以扩张为态射
\[
 \Pi_D^{K,n}:\overline M_D^{K+1,n-K-1}
 \longrightarrow\overline M_D^{K,n-K},
\]
使得：(a) $\Pi$ 与自然嵌入 $M_D^{K,n-K}\hookrightarrow M_D^{K+1,n-K-1}$ 的复合为恒等映射；(b) $\Pi$ 在边界
\[
 \Gamma_D^{K,n}\mathrel{\mathop:}=
 \overline M_D^{K+1,n-K-1}-M_D^{K+1,n-K-1}
\]
上的限制为根态射（radical morphism）。定义
\[
 \overline M_D=\varprojlim_K\overline M_D^{K,n-K},
 \qquad \overline M=\varprojlim_D\overline M_D.
\]
Drinfeld 证明，部分 Frobenius 及 $\GL(2,\A)$ 自然作用于 $\overline M$。现在可以定义
\[
 \check V=H^2(\overline M,\overline\Q_l)
 \mathrel{\mathop:}=\varinjlim_{K,n,D}
 H^2(\overline M_D^{K,n-K},\overline\Q_l).
\]
为定义 $V$，考虑子空间 $V_e\subset\check V$，它由 $\overline M$ 上有理曲线（rational curve）在 $H^2(\overline M,\overline\Q_l)$ 中的类生成。

\paragraph*{注.} 由(3)及 $\overline M$ 上 $\GL(2,\A)$ 作用的存在性可知，每个 $M_D^{K,n-K}$ 上都有无穷多条有理曲线；因此，曲面 $M_D^{K,n-K}$ 含有无穷多条有理曲线，但这些曲线并不组成一个代数族。当 $D$ 很大时，$M_D^{K,n-K}$ 很可能是一般型曲面（surface of general type）。

可以证明，$H^2(\overline M,\overline\Q_l)$ 上典范杯积（cup-product）限制到 $V_e$ 后非退化（nondegenerate）。Drinfeld 将 $V$ 定义为 $V_e$ 在 $\check V$ 中的正交补（orthogonal complement）；也可以将 $V$ 看成商空间 $\check V/V_e$。当然，部分 Frobenius 保持 $V_e$，于是如前所述，在 $V$ 上得到了 $\GL(2,\A)\times W_K\times W_K$ 的表示 $\rho$。

\begin{originaltheorem}{定理 8（Theorem 8）}
表示 $(\rho,V)$ 满足定理6的条件。
\end{originaltheorem}

假定已经完成 $\overline M$ 的构造，接下来如何证明所得表示 $(\rho,V)$ 实现了互反律？

% 原书354页
唯一已知的方法基于迹公式（trace formula）。有下述一般断言。设 $(\rho,V)$ 为 $\GL(n,\A)\times W_K\times W_K$ 的一个分次表示，满足：

(a) 对于 $\GL(n,\A)$ 上的每个 Schwartz--Bruhat 函数（Schwartz--Bruhat function），$\operatorname{Tr}^{\mathrm{reg}}(f)$ 存在；

(b) 对于 $K$ 的任意两个素点 $\mathfrak p',\mathfrak p''$、任意 $m',m''\in\Z^+$ 及 $\GL(n,\A^{\mathfrak p',\mathfrak p''})$ 上的任意 Schwartz--Bruhat 函数 $f$，有等式
\begin{equation}
\begin{split}
 &\operatorname{Tr}^{\mathrm{reg}}_V\bigl((X_{\mathfrak p'}^0\otimes
 X_{\mathfrak p''}^0\otimes f)\times
 \operatorname{Fr}_{\mathfrak p'}^{m'}\times
 \operatorname{Fr}_{\mathfrak p''}^{m''}\bigr)\\
 &\hspace{15mm}=\operatorname{Tr}^{\mathrm{reg}}_L
 \bigl(X_{\mathfrak p'}^{m'}\times\check X_{\mathfrak p''}^{m''}\times f\bigr),
 \qquad \check X_{\mathfrak p}^m(g)\mathrel{\mathop:}=X_{\mathfrak p}^m(g^{-1}).
\end{split}\tag{!}
\end{equation}
则表示 $(\rho,V)$ 实现互反律。

\paragraph*{注.} $\operatorname{Fr}_{\mathfrak p'}^{m'}$ 表示 $W_K$ 上的一个测度（measure），定义如下。先考虑局部群 $W_{K_{\mathfrak p'}}$。它包含惯性群 $I_{\mathfrak p'}$，且 $W_{K_{\mathfrak p'}}/I_{\mathfrak p'}\simeq\Z$。用 $\mu_{\mathfrak p'}^{m'}$ 表示 $W_{K_{\mathfrak p'}}$ 上集中于 $m'$ 的原像、满足
\[
 \int_{W_{K_{\mathfrak p'}}}\mu_{\mathfrak p'}^{m'}=1
\]
的 $I_{\mathfrak p'}$ 不变测度。有一个嵌入 $W_{K_{\mathfrak p'}}\hookrightarrow W_K$（确定到共轭）；用 $\operatorname{Fr}_{\mathfrak p'}^{m'}$ 表示 $\mu_{\mathfrak p'}^{m'}$ 在此嵌入下的像。

据此可用标准方法证明上述断言，证明从略。

为将其用于我们的情形，先要在 $V$ 上给出分次结构，再计算等式两边。分次结构容易得到。代数空间 $M_D$ 是 $M_D^n$ 的不交并。$\overline M_D$ 及 $\overline M=\bigcup\overline M^n$ 也如此。因此可写
\[
 V=\bigoplus V^t,\qquad
 V^t=H^2(\overline M^{2t}\cup\overline M^{2t+1},\overline\Q_l).
\]
有了 $\overline M$ 的定义后，就会看到它满足条件
$\rho(g)V^t=V^{t+v(\det g)}$，$g\in\GL(2,\A)$。

为求右边的 $\operatorname{Tr}^{\mathrm{reg}}(X_{\mathfrak p'}^{m'}\times\check X_{\mathfrak p''}^{m''}\times f)$，可应用 Selberg 迹公式（Selberg trace formula）[7]；幸运的是，此公式对 $\GL(2,\A)$ 已知。

为此，除了相当复杂的计算之外，只需求 $X_{\mathfrak p}^m$ 的“轨道积分”（orbital integral），即对于 $\GL(n,K_{\mathfrak p})$ 的每个共轭类 $\Omega$，求积分 $\int_\Omega X_{\mathfrak p}^m(w)\,dw$。其答案是熟知的（见[8]）。事实上，Drinfeld 已经解决了 $\GL(n,K_{\mathfrak p})$ 的相应问题。下面给出他的结果，它也可用于研究 $\GL(n)$ 的互反律。

\begin{originaltheorem}{定理 9（Theorem 9）}
(1) 若 $\Omega$ 为 $\GL(n,K_{\mathfrak p})$ 中的非半单共轭类，则
\[
 \int_\Omega X_{\mathfrak p}^m\,dw=0.
\]
设 $\Omega$ 为半单共轭类，$\gamma\in\Omega$。则
\[
 Z_G(\gamma)=\prod_i\GL(n_i,L_i),
 \qquad \sum_i n_i[L_i:K_{\mathfrak p}]=n,
\]
其中 $L_i\supset K_{\mathfrak p}$ 为有限扩张。因此可写 $\gamma=\prod\gamma_i$，$\gamma_i\in L_i^*$。

(2) 若 $\int_\Omega X_{\mathfrak p}^m(w)\,dw\ne0$，则存在 $j$ 使得
\[
 n_jv\bigl(N_{L_j/K_{\mathfrak p}}(\gamma_j)\bigr)=m,
 \qquad v\bigl(N_{L_i/K_{\mathfrak p}}(\gamma_i)\bigr)=0\quad(i\ne j).
\]
定义 $\int_\Omega$ 时须在 $\Omega$ 上固定 Haar 测度（Haar measure）。由于 $\Omega=Z_G(\gamma)\backslash\GL(n,K_{\mathfrak p})$，只须对所有 $r\in\Z^*$ 及有限扩张 $L\supset K_{\mathfrak p}$，在 $\GL(r,L)$ 上固定 Haar 测度 $d\gamma$。我们用条件 $\operatorname{meas}\GL(r,o_L)=1$ 来固定它。

(3) 设半单元素 $\gamma\in\Omega_\gamma$ 满足(2)的条件。则
\[
 \int_{\Omega_\gamma}X_{\mathfrak p}^m(w)\,dw
 =(1-q_i)(1-q_i^2)\cdots(1-q_i^{n_i-1})[l_i,k],
\]
其中 $l_i$ 为 $L_i$ 的剩余域，$q_i=\operatorname{card}l_i$。
\end{originaltheorem}

% 原书355页
这一定理的证明基于 Drinfeld 找到的 $X_{\mathfrak p}^m$ 的显式公式。他证明：

(1) $X_{\mathfrak p}^m(g)\ne0$ 蕴含 $g\in\GL(n,o_{\mathfrak p})$ 且 $v(\det g)=m$；\footnote{\textbf{译者注（本次校订）：}这里的整性条件应为 $g\in M_n(o_{\mathfrak p})\cap\GL(n,K_{\mathfrak p})$。若 $g\in\GL(n,o_{\mathfrak p})$，其行列式是单位，必有 $v(\det g)=0$，与本文取 $m>0$ 矛盾。下文讨论约化矩阵可能降秩，也要求采用前一个条件。}

(2) 若 $g$ 满足(1)，则
\[
 X_{\mathfrak p}^m(g)=(1-q)(1-q^2)\cdots(1-q^{r-1}),
\]
其中 $n-r$ 是 $g$ 模 $\pi$ 的约化矩阵的秩。

最后计算(!)的左边。只须考虑 $f$ 为集合 $U_DgU_D$ 的示性函数（characteristic function）的情形，其中 $g\in G$，$D$ 为 $X$ 的除子，$U_D\subset\GL(n,\A)$ 为模 $D$ 等于恒等矩阵的整阿代尔矩阵构成的同余紧开子群（congruence compact open subgroup）。为此固定某个 $N\geq1$，令
\[
 {}^N\overline M_D=
 \bigcup_n\overline M_D^{[(n+N)/2],[(n-N+1)/2]}.
\]
空间 ${}^N\overline M_D$ 是有限型代数空间的分次并。由定义，有投影 $\Pi:\overline M\longrightarrow{}^N\overline M_D$。用
$\Gamma_{g,D}\subset{}^N\overline M_D\times{}^N\overline M_D$ 表示映射
\[
 \Pi\times\Pi\circ g:\overline M
 \longrightarrow{}^N\overline M_D\times{}^N\overline M_D
\]
的像。此 $\Gamma_{g,D}$ 定义了 ${}^N\overline M_D$ 上的一个对应（correspondence）。Drinfeld 证明，${}^N\overline M_D$ 是拟光滑曲面（quasi-smooth surface），即以任意几何点为支集的平展上同调（étale cohomology）与光滑情形相同。因此 $\Gamma_{g,D}$ 在 $H^2({}^N\overline M_D,\overline\Q_l)$ 上定义了一个算子 $A_{g,D}$。

簇 ${}^N\overline M_D$ 是 $X\times X$ 的泛点 $\eta$ 上的曲面。对几乎所有闭点（closed point）$(\mathfrak p',\mathfrak p'')$，它在该点有良约化（good reduction）$R$。$R$ 的上同调与 $H^*({}^N\overline M_D,\overline\Q_l)$ 同构，可在 $H^2(R,\overline\Q_l)$ 上考虑算子
\[
 C_{g,D}=\operatorname{Fr}_{\mathfrak p'}^{m'}\otimes
 \operatorname{Fr}_{\mathfrak p''}^{m''}\otimes A_{g,D}.
\]
(!)左边归结为计算 $\operatorname{Tr}^{\mathrm{reg}}C_{g,D}$。为此应用 Lefschetz 迹公式（Lefschetz trace formula，见[8]），得到\footnote{\textbf{译者注（本次校订）：}Lefschetz 迹公式左端应从 $i=0$ 开始求和，包含 $H^0$ 项；这也与紧接下文同时计算 $H^0$ 和 $H^4$ 的做法一致。此外，右端应采用诱导算子 $C_{g,D}$ 的完整对应，即将 $\Gamma_{g,D}$ 与两侧的 Frobenius 幂复合所得的对应。若只取 $\Gamma_{g,D}$ 的约化，右端便与 $m',m''$ 无关，不能等于一般依赖这两个指数的左端。}
\[
 \sum_{i=1}^{\infty}(-1)^i\Tr\bigl(C_{g,D}\mid
 H^i({}^N\overline M_D,\overline\Q_l)\bigr)
 =\sum_{x\in\Delta\cap\overline\Gamma_{g,D}}n(x),
\]
其中 $\Delta\subset R\times R$ 为对角线，$\overline\Gamma_{g,D}$ 为 $\Gamma_{g,D}$ 的约化。

先看左边。利用 $\overline M$ 上两个有理曲线系统的存在性，可证明
\[
 H^1({}^N\overline M_D,\overline\Q_l)
 =H^3({}^N\overline M_D,\overline\Q_l)=0.
\]
$H^0$ 和 $H^4$ 容易求得，因此只剩计算右边。

为此，Drinfeld 将 $\overline M$ 在 $(\mathfrak p',\mathfrak p'')$ 处约化的几何点描述为
$\operatorname{Fr}_{\mathfrak p'}\times\operatorname{Fr}_{\mathfrak p''}\times\GL(2,\A^{\mathfrak p',\mathfrak p''})$ 模。\footnote{$\A^{\mathfrak p',\mathfrak p''}$ 表示去掉 $\mathfrak p'$ 和 $\mathfrak p''$ 两个分量的阿代尔。}这给出了 $\Delta\cap\overline\Gamma_{g,D}$ 的描述。若 $R$ 光滑，就能断定当 $m',m''$ 很大时，$n(x)=1$。利用 $R$ 拟光滑这一事实，他证明这一结论渐近成立（对很大的 $m',m''$），再与 Selberg 迹公式的相应一边比较，即得(!)。

\begin{thebibliography}{13}
\bibitem[1]{c38-1} P. Cartier, \textit{Representations of a $p$-adic group: A survey}, 本论文集，第1册，111--155.
\bibitem[2]{c38-2} P. Deligne, \textit{Formes modulaires et représentations de $\GL(2)$}, Lecture Notes in Math., vol. 349, Springer-Verlag, Berlin and New York.
\bibitem[3]{c38-3} P. Deligne, \textit{Les constantes des équations fonctionnelles des fonctions $L$}, Lecture Notes in Math., vol. 349, Springer-Verlag, Berlin and New York.
\bibitem[4]{c38-4} P. Deligne, \textit{Cohomologie étale}, SGA $4\frac12$, Lecture Notes in Math., vol. 569, Springer-Verlag, Berlin and New York.
\bibitem[5]{c38-5} D. Flath, \textit{Decomposition of representations into tensor products}, 本论文集，第1册，179--183.
\bibitem[6]{c38-6} I. M. Gelfand, M. I. Graev and I. I. Piatetski-Shapiro, \textit{Representation theory and automorphic functions}, Saunders, 1969.
\bibitem[7]{c38-7} H. Jacquet and R. P. Langlands, \textit{Automorphic forms on $\GL(2)$}, Lecture Notes in Math., vol. 114, Springer-Verlag, Berlin and New York.
\bibitem[8]{c38-8} R. P. Langlands, \textit{Modular forms and $l$-adic representations}, Lecture Notes in Math., vol. 349, Springer-Verlag, Berlin and New York.
\bibitem[9]{c38-9} I. I. Piatetski-Shapiro, \textit{Converse theorem for $\GL(3)$}, Mimeo Notes No. 15, Univ. of Maryland, 1975.
\bibitem[10]{c38-10} I. I. Piatetski-Shapiro, \textit{Multiplicity one theorems}, 本论文集，第1册，209--212.
\bibitem[11]{c38-11} J.-P. Serre, \textit{Représentation linéaires des groupes finis}, Hermann, Paris, 1967.
\bibitem[12]{c38-12} J. Tate, \textit{Number theoretic background}, 本论文集，第2册，3--26.
\bibitem[13]{c38-13} H. Weyl, \textit{The classical groups, their invariants and representations}, Princeton Univ. Press, Princeton, 1939.
\end{thebibliography}
\begin{flushright}Harvard University\end{flushright}
